% ===============================================================
% A4单页项目计划书 LaTeX 模板
% 强烈推荐使用 XeLaTeX 编译器进行编译以获得最佳中文支持
% ===============================================================

\documentclass[a4paper, 12pt]{article}

% --- 1. 导言区 (Preamble) ---
\usepackage{ctex}
\usepackage{graphicx}
\usepackage{caption}
\usepackage[left=2.5cm, right=2.5cm, top=2.5cm, bottom=2.5cm]{geometry}
\usepackage{titlesec}
\titleformat{\section}{\Large\bfseries}{\thesection}{1em}{}[\titlerule]

\newcommand{\projectitem}[1]{\item \textbf{#1}:}

% --- 2. 文档信息 ---
\title{\bfseries 项目计划书}
\author{}
\date{}

% --- 3. 文档正文 ---
\begin{document}

\begin{center}
    {\Huge\bfseries 项目计划书}
    \vspace{1cm}
\end{center}

% --- 项目基本信息 ---
\section*{项目基本信息}

\begin{itemize}
    \projectitem{项目名称}
    夏日海上列车 —— 网页互动叙事游戏

    \projectitem{项目口号}
    乘着海风，驶入一段奇幻旅程

    \projectitem{项目成员}
    \begin{itemize}
        \item 马骏逸 (组长): 负责项目管理、Three.js场景搭建与核心交互开发，以及各系统开发。
        \item 黎佳易 (组员): 负责模型、动画，UI/UX设计与场景美术优化
    \end{itemize}
\end{itemize}

% --- 项目简介 ---
\section*{项目简介}

本项目是一款基于浏览器的3D互动叙事游戏。玩家将置身于一列行驶在无垠海面上的列车中，视角与列车同步向前，沿途邂逅风格各异的海洋生物与神秘角色，触发随机或固定剧情事件，并与之互动。游戏以"旅途"为核心体验，强调沉浸感、发现感与情感共鸣，旨在用轻量的网页技术承载一段富有诗意的微型冒险。

% --- 核心技术点 ---
\section*{核心技术点}

\begin{itemize}
    \projectitem{前端框架与3D引擎}
    采用 \textbf{Three.js} 作为核心渲染引擎，搭配 \textbf{Vite} 构建工具，实现轻量化、热重载的开发流程。

    \projectitem{轨道与列车系统}
    \begin{itemize}
        \item 使用CatmullRom曲线等预设海上轨道，列车沿曲线匀速行进。
        \item 相机绑定列车位置，同时保留可旋转视角让玩家环顾四周。
    \end{itemize}

    \projectitem{海洋渲染}
    \begin{itemize}
        \item \textbf{自定义着色器 (GLSL)}: 基于Voronoi F1 - SmoothF1边缘检测的动漫风格海面，呈现宫崎骏式的色块化波浪。
        \item \textbf{多层波浪系统}: 正弦波叠加 + 顶点位移，模拟柔和起伏的海面动态。
        \item \textbf{距离淡出}: 海面随视距渐变透明，营造无限海洋的视觉感受。
    \end{itemize}

    \projectitem{角色与生物系统}
    \begin{itemize}
        \item 使用 GLTF 格式导入动画角色和生物模型，支持骨骼动画混合。
        \item 生物随列车行进出现，部分生物可与列车同行。
    \end{itemize}

    \projectitem{事件与剧情系统}
    \begin{itemize}
        \item 基于路点触发：列车到达特定坐标时自动触发剧情事件。
        \item 分支对话：玩家选择影响后续剧情走向，支持简单状态机管理剧情状态。
    \end{itemize}

    \projectitem{视觉与特效}
    \begin{itemize}
        \item \textbf{粒子系统}: 用于海浪水花、黄昏萤火等氛围效果。
        \item \textbf{昼夜循环}: 动态切换天空颜色、光照强度与方向，实现从日出到星夜的完整旅程。
        \item \textbf{风格化着色}: 结合色带（ColorRamp）与边缘光效，打造吉卜力风格的视觉语言。
    \end{itemize}

    \projectitem{性能优化}
    \begin{itemize}
        \item 模型使用Draco压缩，控制面数在移动端可接受范围。
        \item InstancedMesh用于重复物体，降低Draw Call。
        \item LOD系统：远处生物使用简化模型或Billboard代替。
    \end{itemize}
\end{itemize}

\vspace{2cm}

\begin{flushright}
    \textbf{计划制定日期:} \today
\end{flushright}

\end{document}
